\section{Resultados Esperados}
\label{sc:resultados}

\subsection{Resultados Preliminares (Iniciação Científica)}
A viabilidade técnica do ecossistema encontra-se consolidada pela IC prévia: (i) pipeline Medallion funcional orquestrado por Dagster, processando mais de 4,05 milhões de registros; (ii) tabela analítica no Apache Druid respondendo consultas OLAP multidimensionais em milissegundos; (iii) grafo topológico do SIN em NetworkX para análise de criticidade $N-1$; e (iv) agrupamento espaço-temporal por ST-DBSCAN validado com sucesso.

\subsection{Resultados Esperados para o TCC02}
Com a execução do plano proposto, espera-se alcançar:
\begin{enumerate}
    \setlength\itemsep{0em}
    \item \textbf{Precisão Geométrica de Risco:} Eliminação dos alarmes falsos a barlavento característicos de buffers circulares estáticos, projetando a mancha de perigo na direção do avanço convectivo do fogo;
    \item \textbf{Aceleração Algorítmica Geoespacial:} Redução drástica da latência de cruzamento espacial em tempo real, viabilizando \textit{hash joins} em tempo constante $O(1)$ com células hexagonais Uber H3 (\texttt{uint64});
    \item \textbf{Despacho Proativo e Relatórios Instantâneos (SitRep):} Redução do tempo de triagem de incidentes de horas para menos de 30 segundos, mediante disparos automatizados de alertas via Webhooks e relatórios PDF no padrão ICS-209 com mapa vetorial anexado;
    \item \textbf{Democratização da Análise Operacional com IA:} Interface conversacional determinística baseada em FastMCP (\textit{Text-to-SQL}), permitindo que operadores consultem o banco colunar em linguagem natural sem alucinações;
    \item \textbf{Publicação Científica e Transferência:} Submissão em conferência qualificada e disponibilização do código aberto e reprodutível.
\end{enumerate}
